\section{Student and Methods}
\label{sec:methods}

\subsection{Student and training loop}
All methods train the same student: brax's PPO
networks, loss and hyperparameters from MuJoCo Playground's locomotion
recipe (policy and value MLPs $512\text{-}256\text{-}128$ with ELU, running
observation normalisation, 20-step unrolls, 4 epochs $\times$ 32
minibatches, $\gamma=0.97$, $\lambda=0.95$, learning rate $3\times10^{-4}$,
entropy cost $10^{-2}$, clipping 0.2). Brax's own training loop fixes each
environment's model for a whole run, so we drive the same components from a
custom loop. One outer \emph{iteration} runs 256 lanes $\times$ 32
environments; every lane holds one level, all 32 of its environments are
reset onto it with a fresh physics-randomisation draw, and the batch is
stepped for 300 steps in 20-step unrolls, with a PPO update after each
unroll when the iteration trains. An iteration is therefore
$256\cdot32\cdot300\approx2.46$M environment steps. Each lane's score is
either the positive value loss, the mean positive GAE advantage over the
iteration~\citep{jiang2021plr}, or the learnability $p(1-p)$, where $p$ is
the fraction of the lane's environments that never fell and travelled at
least 1\,m ($\times s$) from spawn~\citep{rutherford2024noregrets}.
Runs last 300M environment steps for Wojtek and 1B for Go1, with three
seeds per method and robot.

\subsection{Methods}
\label{sec:methodlist}

\paragraph{DR.} Every iteration trains on 256 fresh uniform levels.

\paragraph{DR-easy.} DR on ACCEL's easy generator: a uniform genome shrunk
toward the flat genome by a factor 0.1. It isolates the effect of
``training on easy levels'' from the effect of curation.

\paragraph{\plr~\citep{jiang2021replay}.} A buffer of the 4000
highest-scoring levels (positive value loss). Once it holds at least 256
levels, an iteration replays with probability 0.8: it samples 256 levels by
rank prioritisation (temperature 0.3) mixed with staleness (weight 0.3),
trains on them and refreshes their scores. Otherwise it scores 256 fresh
uniform levels without training and adds them to the buffer.

\paragraph{ACCEL~\citep{parkerholder2022accel}.} \plr\ with two changes:
new levels come from the easy generator instead of the uniform one, and
every replayed level is mutated by perturbing two randomly chosen
parameters by $\mathcal{U}(-0.1,0.1)$ (clipped to $[0,1]$, seed kept). The
next non-training iteration scores the children of the most recent replay
batch (children of earlier consecutive replay batches are discarded), and
fresh easy levels fill the remaining lanes.

\paragraph{ACCEL-SFL.} ACCEL with two changes: the learnability score
$p(1-p)$ in place of positive value loss, and a ``grow'' mutation that,
after the local edit, scales the genome's distance from flat by
$\mathcal{U}(1,1.3)$. Unlike SFL~\citep{rutherford2024noregrets}, which
samples uniformly and keeps a learnability buffer, it keeps ACCEL's
generator and replay. ACCEL and ACCEL-SFL therefore differ in two
switches, score and mutation, and a difference between them cannot be
attributed to the score alone. The pilot (\S\ref{sec:pilot}) ran the
missing arm, ACCEL-grow (grow mutation, positive value loss); the main
study does not.

\paragraph{PAIRED~\citep{dennis2020paired}.} A learned generator maps a
16-dimensional Gaussian latent through an MLP to a sigmoid-squashed
Gaussian over the 8 parameters; a random seed is appended. Each iteration
it proposes 256 levels, on which a protagonist and an antagonist (two
independent brax students) both train. The generator's reward per level is
$\max(0,\ \max R_{\text{ant}} - \overline{R}_{\text{pro}})$, the best
antagonist return minus the mean protagonist return over the lane, and it
is updated by single-step PPO (5 epochs, entropy coefficient 0.1). We do
not use the behaviour-cloning distillation between the students
proposed by~\citet{mediratta2023stabilizing}. The budget
counts protagonist steps only, so PAIRED uses twice the simulation of the
other methods.

\paragraph{Budget accounting.}
All methods get the same number of environment steps. Non-training
iterations of \plr, ACCEL and ACCEL-SFL count toward the budget, so these
methods take fewer gradient steps than DR. We report wall-clock and
simulated steps per method in Table~\ref{tab:compute}.
